\section{Discussion}\label{sec:discussion}

The purpose of this benchmark is not to replace the practical pipeline of the blinded precursor, but to determine which conclusions remain when the temporal split, optimizer budget, validation objective, and economic endpoint are made explicit. The resulting evidence separates a model's average locked-test outcome, its sensitivity to the optimizer--seed block, and its simplified trading consequence. This distinction matters for intraday financial prediction, where a change in selection criterion can alter a model ranking without demonstrating a stable change in market-wide forecasting ability.

\subsection{What the Controlled Comparison Establishes}

The present design preserves the precursor's technical-indicator workflow and fixed seven-feature representation, but changes the comparison unit. Rather than reporting a GA-selected configuration after randomized cross-validation, every backbone is evaluated with RS, GA, PSO, DE, and GWO under the same 1,000-candidate allowance and the same chronological train/validation/test partition. Random Search is therefore a budget-matched reference against which directed population methods can be interpreted \cite{bergstra2012random_search}. The common test segment also prevents final performance values from selecting hyperparameters.

This control improves internal comparability, but it narrows the claim that the data can support. The evidence concerns four classifier families on one five-minute WIN period, conditional on an inherited input representation. It does not show that the seven indicators would be selected again in another training period, nor that an optimizer with a favourable result here will retain that advantage under another market regime. The contribution is consequently a controlled experimental frame, not a universal ordering of metaheuristics.

\subsection{Predictive Performance Depends on the Validation Objective}

The locked-test summaries in Table~\ref{tab:predictive_summary} show a clear reversal at the model level. Under MCC/$F_1$ fitness, RF has the largest marginal means for accuracy, MCC, and $F_1$, and its small standard deviations indicate comparatively low sensitivity to the evaluated optimizer--seed blocks. Under weighted-accuracy fitness, MLP is first on all three reported predictive measures, while 1D-CNN is also stable but remains below MLP in mean accuracy and MCC. The change in ordering is substantive: selecting for weighted-accuracy fitness does not simply rescale the MCC/$F_1$ result.

MCC is useful because it summarizes the association between predicted and observed binary classes beyond raw accuracy \cite{chicco2020mcc}; however, optimizing an imbalance-aware criterion cannot guarantee the best outcome for every backbone and search trajectory. Conversely, the MLP advantage under weighted-accuracy fitness is concentrated in a low-variance regime, whereas MLP, SVM, and 1D-CNN outcomes under MCC/$F_1$ fitness vary much more across the evaluated blocks. A reported mean must therefore be accompanied by its dispersion and by the selection rule that generated it.

The rank analysis reinforces this caution. Under MCC/$F_1$ fitness, RF has the highest marginal mean MCC, but MLP has the lowest mean Friedman rank (1.533) across the matched blocks. This is not a contradiction: means weight score magnitude, whereas ranks summarize ordering within each block. Under weighted-accuracy fitness, both summaries place MLP first, and its Holm-adjusted contrasts against all other backbones are significant. Because all 15 blocks share one locked chronological segment, these results indicate separation within the evaluated design rather than independent replication across time.

\subsection{Optimizer Trajectories Are Diagnostic, Not a Universal League Table}

Figures~\ref{fig:rf_convergence_mccf1}--\ref{fig:cnn_convergence_accuracy} index progress by candidate evaluations rather than by optimizer-specific generation counts. Each line shows the mean trajectory over the three common seeds; variability envelopes are omitted to keep closely spaced optimizer curves legible. In each panel, the main y-axis is tightly bounded by the observed range of the five mean trajectories, with a small margin, and an inset magnifies the final 250 evaluations. Because limits are panel-specific, visual heights should not be compared across panels without reading their y-axis scales. Early gains indicate how quickly an optimizer locates useful regions of a search space, whereas later plateaus show the incremental benefit of additional evaluations under the fixed budget. Convergence in validation fitness does not establish equivalent generalization, so the locked-test values and their dispersion in Table~\ref{tab:predictive_summary} remain necessary for interpretation.

Search behaviour interacts with the backbone and the objective. A method that improves early for RF need not be the most reliable procedure for MLP or 1D-CNN, whose mixed discrete--continuous hyperparameter spaces have different response surfaces. The article therefore does not claim a universal winner among RS, GA, PSO, DE, and GWO. A stronger optimizer-level claim would require a dedicated analysis across more seeds and temporal replications, with the optimizer itself as the inferential factor.

\subsection{Predictive Quality and Economic Outcome Are Distinct Endpoints}

Table~\ref{tab:economic_summary} makes the distinction between predictive evaluation and the exploratory trading rule concrete. Under MCC/$F_1$ fitness, RF combines the largest mean net profit with the smallest average drawdown and the lowest drawdown dispersion. Under weighted-accuracy fitness, MLP has the largest mean net profit, profit factor, and annualized Sharpe ratio, while RF retains the least negative average drawdown. Thus, the preferred choice depends on whether the decision criterion prioritizes return, drawdown, or consistency; it cannot be inferred from accuracy alone.

These values are deliberately reported as points under a fixed intrabar long/short rule with a five-point round-trip cost. They are not an investable-strategy claim: the evaluation does not model a full transaction-cost schedule, market impact, delayed execution, position sizing, or a live signal-selection policy. Its role is diagnostic: it shows that converting a directional label to trades can change the practical ranking of model families, and that economic reporting should accompany predictive metrics.

\subsection{Limitations and Research Agenda}

Four limitations delimit the present evidence. First, only three seeds are common to every model--optimizer cell; standard deviations and rank tests describe the evaluated blocks, not repeated market samples. Second, the locked test covers one instrument and one chronological segment, so regime sensitivity and cross-market transfer remain unknown. Third, the Information-Gain mask was inherited from the blinded precursor rather than re-estimated on the current training segment; all conclusions are conditional on that representation. Fourth, equalizing the budget at 1,000 evaluations establishes fairness at one computational scale but does not characterize larger-budget behaviour or wall-clock efficiency.

The next study should preserve the current leakage-aware split and add a preregistered common 20-seed design across multiple non-overlapping market periods. It should extend the present naive and default-RF references to no-tuning baselines for every backbone; archive every candidate evaluation; report compute time; and test optimizer differences separately from model differences. Re-estimating feature selection within each training window would further distinguish the effect of representation from hyperparameter search. These additions would turn the current transparent benchmark into confirmatory evidence without weakening its temporal protocol.
